\section{Symbols, Notation, and Units}

\subsection{Ctrl+F keywords: symbol, notation, units, constant, subscript, standard state, phase, concentration, activity}
Use this index when a formula contains an unfamiliar quantity or mark. Topic sections repeat the symbols needed for their calculations.

\subsection{Python toolbox reference}
The calculation notebook is \texttt{exam\_toolbox.ipynb}; it imports \texttt{exam\_tools} from \texttt{support\_material/input/exams}. Use it when a recipe below names a Python function, especially for formula parsing, nonlinear equilibrium roots, logarithmic/exponential relations, regression, or repeated unit conversions. The toolbox also exports physical constants; use \texttt{constant(...)} or \texttt{constant\_info(...)} to check values such as \(R\), \(F\), \(N_{\mathrm A}\), \(h\), and \(c\).

\subsection{Common calculation quantities}
\referencetable{tab:symbols:calculation}{Common calculation quantities and typical units.}
\begin{center}
\begin{tabular}{@{}>{$}l<{$} p{7.2cm} l@{}}
\toprule
\text{Symbol} & \text{Meaning} & \text{Typical unit}\\
\midrule
n & Amount of substance; context can change this meaning as listed below & \(\mathrm{mol}\)\\
m & Mass; in colligative-property formulas, italic \(m\) may instead mean molality & \(\mathrm{g}\), \(\mathrm{kg}\), or \(\mathrm{mol\,kg^{-1}}\)\\
M & Molar mass & \(\mathrm{g\,mol^{-1}}\)\\
c,\ [\ce{A}] & Amount concentration of a species \(\ce{A}\) in solution & \(\mathrm{mol\,L^{-1}}\)\\
V & Volume & \(\mathrm{L}\) or \(\mathrm{m^3}\)\\
p,\ P & Pressure or vapor pressure, according to context & \(\mathrm{bar}\), \(\mathrm{atm}\), or \(\mathrm{Pa}\)\\
T & Absolute temperature & \(\mathrm{K}\)\\
\rho & Density & \(\mathrm{g\,cm^{-3}}\) or \(\mathrm{kg\,m^{-3}}\)\\
R & Gas constant, used with units matching the formula data & e.g. \(\mathrm{J\,mol^{-1}K^{-1}}\)\\
N_{\mathrm A} & Avogadro constant, \(6.022\times10^{23}\) particles per mole & \(\mathrm{mol^{-1}}\)\\
\nu_i & Stoichiometric coefficient of species \(i\) & dimensionless\\
\xi & Reaction extent & \(\mathrm{mol}\)\\
\bottomrule
\end{tabular}
\end{center}

\subsection{Thermodynamics, equilibrium, kinetics, and electrochemistry}
\referencetable{tab:symbols:advanced}{Thermodynamics, equilibrium, kinetics, and electrochemistry symbols.}
\begin{center}
\begin{tabular}{@{}>{$}l<{$} p{7.2cm} l@{}}
\toprule
\text{Symbol} & \text{Meaning} & \text{Typical unit}\\
\midrule
q,\ w & Heat transferred to the system; work done on the system & \(\mathrm{J}\) or \(\mathrm{kJ}\)\\
U,\ H,\ S,\ G & Internal energy, enthalpy, entropy, Gibbs free energy & \(\mathrm{J}\), \(\mathrm{kJ}\), or \(\mathrm{J\,K^{-1}}\)\\
\Delta & Change: final value minus initial value, or reaction change & as quantity\\
Q & Reaction quotient computed from current composition & dimensionless\\
K & Equilibrium constant; subscripts identify the equilibrium type & dimensionless\\
r & Reaction rate & \(\mathrm{mol\,L^{-1}s^{-1}}\)\\
k & Rate constant; units depend on reaction order & order-dependent\\
E_{\mathrm a} & Activation energy & \(\mathrm{J\,mol^{-1}}\)\\
E_{\mathrm{cell}} & Electrochemical cell potential & \(\mathrm V\)\\
F & Faraday constant, charge per mole of electrons & \(\mathrm{C\,mol^{-1}}\)\\
\bottomrule
\end{tabular}
\end{center}

\subsection{Marks, subscripts, and chemical notation}
\referencetable{tab:symbols:notation}{Marks, subscripts, phases, and chemical notation.}
\begin{center}
\begin{tabular}{@{}>{$}l<{$} p{10.2cm}@{}}
\toprule
\text{Notation} & \text{Meaning}\\
\midrule
{}^\circ & Standard-state value, normally referring to stated standard conditions\\
\mathrm{rxn},\ \mathrm f,\ \mathrm{comb} & Reaction, formation, and combustion quantity\\
\mathrm{vap},\ \mathrm{fus},\ \mathrm{sub} & Vaporization, fusion, and sublimation\\
\mathrm{tot},\ \mathrm{ext},\ 0 & Total, external, and initial quantities\\
\mathrm{sp},\ \mathrm{cell},\ \mathrm{eq} & Solubility product, electrochemical cell, and equilibrium\\
a_{\ce{A}} & Activity of species \(\ce{A}\); dilute-solution calculations often use concentration in its place\\
(s),(l),(g),(aq) & Solid, liquid, gas, and aqueous phase\\
\ce{->},\ \ce{<=>} & Reaction arrow and reversible/equilibrium arrow\\
\ln,\ \log & Natural logarithm and base-10 logarithm\\
\Sigma,\ \sum & Sum over the indicated species or terms\\
\bottomrule
\end{tabular}
\end{center}

\subsection{Symbols with multiple meanings}
\referencetable{tab:symbols:overloaded}{Symbols whose meaning must be identified from context.}
\begin{center}
\begin{tabular}{@{}>{$}l<{$} p{10.2cm}@{}}
\toprule
\text{Symbol} & \text{Meaning determined by context}\\
\midrule
n & Amount in moles in stoichiometry/gases; principal quantum number in electronic structure; moles of electrons in electrochemistry; degree of polymerization in polymers.\\
M & Molar mass in calculations; the generic metal symbol in dissolution formulas such as \(\mathrm{MX_2}\).\\
Q,\ q & Capital \(Q\) is reaction quotient; lowercase \(q\) is heat.\\
p,\ P & Lowercase \(p\) is generally gas/partial pressure; uppercase \(P\) is used for vapor pressure in the phase section.\\
\bottomrule
\end{tabular}
\end{center}
